\documentclass[a4paper,10pt]{article}
\usepackage[T2A]{fontenc}
\usepackage[utf8]{inputenc}
\usepackage[russian]{babel}
\usepackage{amsmath,amssymb,mathtools}
\usepackage[a4paper,margin=20mm,headheight=22pt,headsep=5mm,footskip=11mm]{geometry}
\usepackage{microtype}
\usepackage{tabularx,booktabs,array}
\usepackage{enumitem}
\usepackage{xcolor}
\usepackage{tikz}
\usepackage{fancyhdr}
\usepackage{hyperref}
\usepackage{parskip}
\definecolor{accent}{HTML}{008080}
\definecolor{darkaccent}{HTML}{005959}
\definecolor{ink}{HTML}{333333}
\definecolor{mistake}{HTML}{B03030}
\hypersetup{colorlinks=true,linkcolor=darkaccent,urlcolor=darkaccent,pdftitle={Ревью домашней работы — Ксения Васильченко},pdfauthor={Лёвин Артём Александрович}}
\setlength{\parindent}{0pt}
\setlength{\parskip}{3pt}
\setlist[enumerate]{topsep=2pt,itemsep=1.5pt,leftmargin=18pt}
\pagestyle{fancy}
\fancyhf{}
\fancyhead[L]{\footnotesize Проценты и скидки}
\fancyhead[R]{\footnotesize Ксения Васильченко}
\fancyfoot[C]{\footnotesize \thepage}
\renewcommand{\headrulewidth}{0.3pt}
\renewcommand{\footrulewidth}{0pt}
\renewcommand{\headrule}{\hbox to\headwidth{\color{accent}\leaders\hrule height \headrulewidth\hfill}}
\newcommand{\sectionline}{\par\vspace{2pt}{\color{accent}\hrule height 0.45pt}\vspace{6pt}}
\newcommand{\tasktitle}[2]{\phantomsection\label{task:#1}{\LARGE\color{darkaccent}\bfseries Задание № #1}\par\vspace{1mm}{\color{ink}#2}\sectionline}
\newcommand{\field}[1]{\par\vspace{2.5pt}{\bfseries\color{darkaccent}#1}\par\vspace{0pt}}
\newcommand{\redmark}{\textcolor{mistake}{\textbf{Ошибка:}}}
\newcommand{\idea}{\textcolor{darkaccent}{\textbf{Ключевая идея:}}}
\newcommand{\outcome}{\textcolor{darkaccent}{\textbf{Итог:}}}
\begin{document}
\hypersetup{pageanchor=false}
\thispagestyle{empty}
\begin{titlepage}
\centering
\vspace*{25mm}
{\normalsize\color{darkaccent} МАТЕМАТИКА\par}
\vspace{21mm}
{\LARGE Ревью домашней работы\par}
\vspace{12mm}
{\Large\color{accent} Проценты. Доли и скидки\par}
\vspace{9mm}
{\normalsize Разбор решений и отработка ошибок\par}
\vfill
{\normalsize Ксения Васильченко\par}
\vspace{10mm}
{\normalsize 14 ноября 2026 г.\par}
\vspace{14mm}
{\small Лёвин Артём Александрович\par}
{\small эксклюзивно для Ксении Васильченко\par}
\vspace{20mm}
\begin{tikzpicture}[x=1cm,y=1cm]
\draw[accent,line width=0.65pt] (0,0) .. controls (1.1,0.55) and (2.5,-0.4) .. (3.9,0.12) .. controls (5.3,0.57) and (6.1,-0.4) .. (7.5,0.15);
\end{tikzpicture}
\end{titlepage}
\hypersetup{pageanchor=true}

\newpage
{\Large\color{darkaccent}\bfseries Сводная таблица}\sectionline
В таблицу включены только задания, в которых обнаружена ошибка в записи вычислений, неправильный результат или недостающая часть ответа. Номера ведут к подробному разбору.
\par\vspace{6pt}
{\small
\renewcommand{\arraystretch}{1.35}
\begin{tabularx}{\linewidth}{@{}p{12mm}p{23mm}X p{25mm}p{28mm}@{}}
\toprule
\textbf{№} & \textbf{Тема} & \textbf{Суть ошибки} & \textbf{Ваш ответ} & \textbf{Правильный ответ} \\
\midrule
\hyperref[task:6]{6} & Скидка & В вычислительной строке записано лишнее деление; числовой результат при этом верен. & 2040 руб. & 2040 руб. \\
\hyperref[task:8]{8} & Доля от целого & Неверно разделены 64 книги на 160: получено 2,5 вместо 0,4. & 250\% & 40\% \\
\hyperref[task:10]{10} & Обратный процент & Найдена первоначальная цена, но не вычислена сумма скидки. & 2000 руб. & 2000 руб.; скидка 300 руб. \\
\bottomrule
\end{tabularx}
}
\par\vspace{12pt}
{\color{darkaccent}\bfseries Основной ориентир}

При нахождении процента от числа умножайте число на соответствующую десятичную дробь. Если известна цена \emph{после} скидки, восстановите исходную цену делением на долю оставшейся стоимости. В ответе обязательно укажите \emph{все} величины, которые требует условие.

\newpage
\thispagestyle{empty}
\vspace*{0.33\textheight}
\begin{center}
{\Huge\color{darkaccent} Разбор ошибок}\par
\vspace{13mm}
{\color{accent}\rule{0.65\textwidth}{0.55pt}}\par
\vspace{8mm}
{\large Задания № 6, 8 и 10}
\end{center}
\newpage

\tasktitle{6}{Нахождение цены после скидки}
\field{Текст задания}
Куртка стоила 2400 рублей. После скидки 15\% сколько рублей нужно заплатить за куртку?
\field{Ваше решение}
Сначала Вы правильно установили, что после скидки остаётся
\[
100\%-15\%=85\%.
\]
Затем в тетради записана строка, содержащая фрагмент \(85:2400\cdot0{,}85=2040\). Само число \(2040\) вычислено верно, но записанная перед множителем операция деления делает равенство некорректным.
\field{Ваш ответ}
2040 рублей.
\field{Ошибка}
\redmark\ некорректная запись арифметического действия при правильном числовом ответе.
\field{Где именно возникает ошибка}
\textbf{Последний правильный шаг:} \(100\%-15\%=85\%=0{,}85\).

\textbf{Первый неверный шаг:} запись с началом \(85:2400\cdot0{,}85\). Если двоеточие обозначает деление, выражение не равно 2040. Запись не позволяет считать это равенство верным.
\field{Почему это неверно}
Величина \(85\%\) обозначает долю от \emph{исходной} цены, а не отдельное число, на которое нужно делить цену. Для нахождения 85\% от 2400 следует \emph{умножить} 2400 на \(0{,}85\). Дополнительная операция деления на 2400 не имеет отношения к условию и меняет значение выражения. Даже если результат в конце строки правильный, каждое записанное равенство должно быть математически обосновано.
\field{Ключевая идея}
\idea\ после уменьшения исходной стоимости на 15\% остаётся 85\% первоначальной цены. Найти эту долю можно одним умножением.
\field{Исправление от последнего правильного шага}
Сохраняем найденные \(85\%\), но убираем лишнее действие:
\[
85\%=0{,}85,\qquad 2400\cdot0{,}85=2040.
\]
\field{Полное правильное решение}
\begin{enumerate}
\item За 100\% принимаем стоимость до скидки: 2400 рублей.
\item После снижения на 15\% осталось \(100\%-15\%=85\%\) исходной суммы.
\item Переводим процент в десятичную дробь: \(85:100=0{,}85\).
\item Находим 85\% от 2400: \(2400\cdot0{,}85=2040\) рублей.
\end{enumerate}
\field{Проверка}
Сама скидка составляет \(2400\cdot0{,}15=360\) рублей. Проверяем вычитанием: \(2400-360=2040\) рублей. Оба способа дают один результат.
\field{Итог}
\outcome\ заплатить нужно \(\boxed{2040\text{ руб.}}\). Ответ верный; нужно исправить запись действия.
\field{Задача на закрепление}
Ботинки стоили 3200 рублей. После скидки 12\% сколько рублей необходимо заплатить? Запишите расчёт без лишних арифметических операций.

\newpage
\tasktitle{8}{Нахождение процента, который часть составляет от целого}
\field{Текст задания}
На выставке 160 книг, из них 64 --- научно-популярные. Какой процент книг составляют научно-популярные?
\field{Ваше решение}
В работе записано:
\[
\frac{64}{160}=2{,}5,\qquad 2{,}5\cdot100=250\%.
\]
\field{Ваш ответ}
250\%.
\field{Ошибка}
\redmark\ неверное деление при определении доли части от целого.
\field{Где именно возникает ошибка}
\textbf{Последний правильный шаг:} дробь \(\frac{64}{160}\) составлена верно: в числителе находится число научно-популярных книг, а в знаменателе --- общее число книг.

\textbf{Первый неверный шаг:} \(\frac{64}{160}=2{,}5\). В действительности \(2{,}5=\frac{160}{64}\); это результат \emph{обратного} деления.
\field{Почему это неверно}
Часть не может оказаться больше целого, если речь идёт о количестве книг одной выставки: 64 книги меньше 160 книг. Значит, отношение \(\frac{64}{160}\) обязано быть меньше единицы, а искомый процент --- меньше 100\%. Число 2,5 больше единицы, а 250\% больше 100\%, поэтому уже простая оценка величины показывает ошибку. Важно не переставлять числитель и знаменатель местами: вопрос «какую часть от 160 составляют 64?» означает именно \(64:160\).
\field{Ключевая идея}
\idea\ чтобы узнать, сколько процентов \emph{часть} составляет от \emph{целого}, разделите часть на целое, затем умножьте результат на 100\%.
\field{Исправление от последнего правильного шага}
Правильно составленную дробь сохраняем и пересчитываем её:
\[
\frac{64}{160}=\frac{4}{10}=0{,}4.
\]
\field{Полное правильное решение}
\begin{enumerate}
\item Все 160 книг принимаем за 100\%.
\item Научно-популярные книги --- это 64 книги из 160, поэтому доля равна \(\frac{64}{160}\).
\item Сокращаем дробь на 16: \(\frac{64}{160}=\frac{4}{10}=0{,}4\).
\item Переводим долю в проценты: \(0{,}4\cdot100\%=40\%\).
\end{enumerate}
\field{Проверка}
Вычислим 40\% от 160: \(160\cdot0{,}4=64\) книги. Получили ровно указанное в условии количество. Кроме того, 40\% меньше 100\%, как и должно быть для части от общего количества.
\field{Итог}
\outcome\ научно-популярные книги составляют \(\boxed{40\%}\) всех книг.
\field{Задача на закрепление}
В библиотеке 180 новых книг, из которых 54 --- исторические. Какой процент новых книг составляют исторические?

\newpage
\tasktitle{10}{Восстановление исходной цены и вычисление суммы скидки}
\field{Текст задания}
После скидки 15\% конструктор стал стоить 1700 рублей. Найдите его первоначальную цену \emph{и величину скидки в рублях}.
\field{Ваше решение}
В тетради Вы записали правильное вычисление первоначальной цены:
\[
1700:\frac{85}{100}=2000.
\]
\field{Ваш ответ}
2000 рублей. Сумма скидки в рублях не указана.
\field{Ошибка}
\redmark\ не завершено решение: найдено только одно из двух запрошенных значений.
\field{Где именно возникает ошибка}
\textbf{Последний правильный шаг:} получена исходная цена 2000 рублей. Деление на \(0{,}85\) выполнено верно.

\textbf{Первый неполный шаг:} после нахождения 2000 рублей решение завершается без вычисления, на сколько рублей уменьшилась цена.
\field{Почему это неверно}
Первоначальная цена и величина скидки --- разные величины. \(2000\) рублей показывает, сколько конструктор стоил \emph{до} снижения, а величина скидки показывает, сколько рублей \emph{вычли} из первоначальной цены. Условие просит оба значения; даже правильно найденная исходная цена не даёт полного ответа. Нельзя также брать 15\% от 1700 рублей: процент скидки отсчитывается от первоначальной, а не от уже уменьшенной цены.
\field{Ключевая идея}
\idea\ после скидки 15\% оплачивают 85\% исходной цены. Сначала находят 100\% по известным 85\%, затем вычисляют разность между исходной и новой ценами.
\field{Исправление от последнего правильного шага}
Сохраняем полученную первоначальную цену 2000 рублей и находим недостающую величину:
\[
2000-1700=300\text{ рублей}.
\]
\field{Полное правильное решение}
\begin{enumerate}
\item До снижения стоимость составляла 100\%, после скидки осталось \(100\%-15\%=85\%\).
\item Получается, что 1700 рублей --- это 85\% исходной стоимости.
\item Переводим 85\% в десятичную дробь: \(85\%=0{,}85\).
\item Находим первоначальную цену: \(1700:0{,}85=2000\) рублей.
\item Находим величину скидки: \(2000-1700=300\) рублей.
\end{enumerate}
\field{Проверка}
Пятнадцать процентов от 2000 рублей --- это \(2000\cdot0{,}15=300\) рублей. Если уменьшить исходную цену на 300 рублей, получим \(2000-300=1700\) рублей, то есть условие выполняется.
\field{Итог}
\outcome\ первоначальная цена --- \(\boxed{2000\text{ руб.}}\), величина скидки --- \(\boxed{300\text{ руб.}}\).
\field{Задача на закрепление}
После скидки 20\% рюкзак стал стоить 1440 рублей. Найдите его первоначальную цену и сумму скидки в рублях.
\end{document}